% @order: 4
% @description: 设置彩色文本框样式

% !文本框设置
% #region

% 推导紫*
\definecolor{PurpleDeep}{HTML}{82218B}       % 核心紫
\definecolor{PurpleLight}{HTML}{F9F2F9}      % 极浅紫背景

% 重点橙
\definecolor{OrangeDeep}{HTML}{FE6304}       % 核心橙
\definecolor{OrangeLight}{HTML}{FFF6F0}      % 极浅橙背景

% 解答绿*
\definecolor{GreenDeep}{HTML}{007355}        % 核心绿
\definecolor{GreenLight}{HTML}{F0F8F6}       % 极浅绿背景

% 禁止红*
\definecolor{RedDeep}{HTML}{D20000}          % 核心红
\definecolor{RedLight}{HTML}{FEF2F2}         % 极浅红背景

% 定义黄
\definecolor{YellowDeep}{HTML}{F59A23}       % 核心黄
\definecolor{YellowLight}{HTML}{FEF9F2}      % 极浅黄背景

% 定理蓝*
\definecolor{BlueDeep}{HTML}{3070B6}         % 核心蓝
\definecolor{BlueLight}{HTML}{F2F6FA}        % 极浅蓝背景

% 例题灰*
\definecolor{GrayDeep}{HTML}{4C5E75}         % 核心灰
\definecolor{GrayLight}{HTML}{F4F6F8}        % 极浅灰背景

% 定义通用样式模板
\tcbset{
    myThemStyle/.style={
        enhanced,
        breakable,
        arc=4pt,            % 圆角大小 
        outer arc=4pt,
        boxrule=0.5pt,      % 普通边框线宽 
        leftrule=3pt,       % 关键：用稍粗的左边框替代 leftbar 参数 
        toprule=0.5pt,
        rightrule=0.5pt,
        bottomrule=0.5pt,
        fonttitle=\bfseries,
        coltitle=black,
        attach title to upper, % 标题嵌入正文
        after title={.\ },     % 标题后的标点
        left=8pt,              % 左边距 
    }
}
\tcbset{
    myTitleStyle/.style={
        enhanced,
        breakable,
        arc=4pt,            % 圆角大小
        outer arc=4pt,
        boxrule=0.5pt,      % 普通边框线宽
        leftrule=3pt,       % 保持左边框加粗的视觉识别度
        toprule=0.5pt,
        rightrule=0.5pt,
        bottomrule=0.5pt,
        fonttitle=\bfseries,
        coltitle=white,     % 标题栏文字默认为白色
        toptitle=3pt,       % 标题栏内部上边距 (微调以垂直居中)
        bottomtitle=3pt,    % 标题栏内部下边距
        left=8pt,           % 内容区左边距
        right=8pt,          % 内容区右边距
        top=4pt,            % 内容区上边距
        bottom=4pt,         % 内容区下边距
    }
}

% 定义“定理”环境
\newcounter{theorem}[section]
\renewcommand{\thetheorem}{\thesection-\arabic{theorem}}
\newtcolorbox{theorem}[1][]{
    myThemStyle,
    colback=BlueLight,   % 背景色
    colframe=BlueDeep,   % 边框色
    before upper={
        \refstepcounter{theorem}
        \textcolor{BlueDeep}{\textbf{定理~\thetheorem}}~
    },
    overlay unbroken and last={
        \node[anchor=south east, inner sep=6pt] 
            at (frame.south east) {\color{BlueDeep}\faAtom};
    },
    #1
}

% 定义“引理”环境
\newcounter{lemma}[theorem]
\renewcommand{\thelemma}{\thesection-(\arabic{lemma})}
\newtcolorbox{lemma}[1][]{
    myThemStyle,
    colback=BlueLight,   % 背景色
    colframe=BlueDeep,   % 边框色
    before upper={
        \refstepcounter{lemma}
        \textcolor{BlueDeep}{\textbf{引理~\thelemma}}~
    },
    overlay unbroken and last={
        \node[anchor=south east, inner sep=6pt] 
            at (frame.south east) {\color{BlueDeep}\faAtom};
    },
    #1
}

% 定义“定义”环境
\newcounter{definition}[section]
\renewcommand{\thedefinition}{\thesection-\arabic{definition}}
\newtcolorbox{definition}[1][]{
    myThemStyle,
    colback=YellowLight,   % 背景色
    colframe=YellowDeep,   % 边框色
    before upper={
        \refstepcounter{definition}
        \textcolor{YellowDeep}{\textbf{定义~\thedefinition}}~
    },
    overlay unbroken and last={
        \node[anchor=south east, inner sep=6pt] 
            at (frame.south east) {\color{YellowDeep}\faLayerGroup};
    },
    #1
}

% 定义“命题”环境
\newcounter{proposition}[subsection]
\renewcommand{\theproposition}{\arabic{proposition}}
\newtcolorbox{proposition}[1][]{
    myThemStyle,
    colback=OrangeLight,   % 背景色
    colframe=OrangeDeep,   % 边框色
    before upper={
        \refstepcounter{proposition}
        \textcolor{OrangeDeep}{\textbf{命题~\theproposition}}~
    },
    overlay unbroken and last={
        \node[anchor=south east, inner sep=6pt] 
            at (frame.south east) {\color{OrangeDeep}\faPencil*};
    },
    #1
}

% 定义“推导”环境
\newtcolorbox{derivation}[1][]{
    myThemStyle,
    colback=PurpleLight,   % 背景色
    colframe=PurpleDeep,   % 边框色
    before upper={
        \textcolor{PurpleDeep}{\textbf{推导}}~
    },
    overlay unbroken and last={
        \node[anchor=south east, inner sep=6pt] 
            at (frame.south east) {\color{PurpleDeep}\faRandom};
    },
    #1
}

% 定义“误区”环境
\newtcolorbox{mistake}[1][]{
    myTitleStyle,
    colback=RedLight,   % 背景色
    colframe=RedDeep,   % 边框色
    before upper={
        \textcolor{RedDeep}{\textbf{误区}}~
    },
    overlay unbroken and last={
        \node[anchor=south east, inner sep=6pt] 
            at (frame.south east) {\color{RedDeep}\faExclamation};
    },
    #1
}

% 定义“问题”与“解答(证明)”环境
\newcounter{question}[subsection]
\renewcommand{\thequestion}{\arabic{question}}
\NewTColorBox{question}{o}
{
    myThemStyle,
    colback=GrayLight,
    colframe=GrayDeep,
    before upper={
        \IfValueTF{#1}
        {
            \textcolor{GrayDeep}{\textbf{问题~#1}}~
        }
        {
            \refstepcounter{question}
            \textcolor{GrayDeep}{\textbf{问题~\thequestion}}~
        }
    },
    overlay unbroken and last={
        \node[anchor=south east, inner sep=6pt] 
            at (frame.south east) {\color{GrayDeep}\faQuestion};
    },
}
\NewTColorBox{solution}{o}
{
    myTitleStyle,
    colback=GreenLight,
    colframe=GreenDeep,
    colbacktitle=GreenDeep,
    title={\IfValueTF{#1}{#1\enspace 解答}{解答}},
    after title={\hfill \faAngleDown},
    overlay unbroken and last={
        \node[anchor=south east, inner sep=6pt] 
            at (frame.south east) {\color{GreenDeep}\faCheckDouble};
    },
}
% 定义“证明”环境
\makeatletter
\@ifundefined{proof}{}{\let\proof\relax \let\endproof\relax}
\makeatother
\NewTColorBox{proof}{o}
{
    myTitleStyle,
    colback=GreenLight,
    colframe=GreenDeep,
    colbacktitle=GreenDeep,
    title={\IfValueTF{#1}{#1\enspace 证明}{证明}},
    after title={\hfill \faAngleDown},
    overlay unbroken and last={
        \node[anchor=south east, inner sep=6pt] 
            at (frame.south east) {\color{GreenDeep}\faSquare[regular]};
    },
}

\newtcolorbox{extension}[2][]{
    myTitleStyle,               % 调用通用标题栏样式
    colback=PurpleLight,         % 内容背景色（建议根据需要调整色值）
    colframe=PurpleDeep,         % 边框颜色
    colbacktitle=PurpleDeep,     % 标题栏背景色
    title={#2},                 % 核心修改：使用第二个必选参数作为标题
    after title={\hfill \faAngleDown},
    overlay unbroken and last={
        \node[anchor=south east, inner sep=6pt] 
            at (frame.south east) {\color{PurpleDeep}\faPaperclip};
    },
    #1                          % 允许传入额外参数（如调整边距等）
}
% #endregion